\begin{explanationbox}
\textbf{\textcolor{CExplanation}{先看具体的倒数对}}
首页曲线上，$t=1/2$ 与 $t=2$ 都给出 $u=5/2$。因此用 $u$ 可以合并两个原变量，但不能再区分它们。
\begin{center}
\begin{tikzpicture}[x=1cm,y=0.62cm,font=\small]
\node[font=\bfseries] at (5.7,5.3) {同一倒数和，差式却反号};
\draw[-Stealth] (1,0)--(10.5,0);
\node[anchor=east] at (1,0) {$0$};
\draw[CStatement,very thick] (3,0)--(3,2.5);
\draw[CStatement,very thick] (8,0)--(8,2.5);
\fill[CStatement] (3,2.5) circle (2pt);\fill[CStatement] (8,2.5) circle (2pt);
\node[anchor=south] at (3,2.8) {$u=5/2$};\node[anchor=south east] at (8,2.8) {$u=5/2$};
\draw[CTrap,very thick,dashed] (3.55,0)--(3.55,-3.75);
\draw[CExample,very thick,dashed] (8.55,0)--(8.55,3.75);
\node[anchor=north,text=CTrap] at (3.55,-4.0) {$t^2-t^{-2}=-15/4$};
\node[anchor=south,text=CExample] at (8.55,4.0) {$t^2-t^{-2}=15/4$};
\node[anchor=north east] at (3.3,-0.3) {$t=1/2$};\node[anchor=north] at (8,-0.3) {$t=2$};
\end{tikzpicture}
\end{center}

实线段表示倒数和，虚线段表示差式，长度按实际数值绘制。互换 $t$ 与 $1/t$ 时，$t^n+t^{-n}$ 不变；$t^n-t^{-n}$ 反号。因此
\[
(t^2-t^{-2})^2=u^2(u^2-4),
\]
但不能不加条件地写成 $t^2-t^{-2}=u\sqrt{u^2-4}$。平方后的式子可以整体换元，差式本身还须判断符号。

\medskip
\textbf{\textcolor{CExplanation}{为什么交叉项容易处理}}
因为 $t\cdot t^{-1}=1$，平方时的两个交叉项固定为 $2$；相乘得到更高次时，中间项又回到较低次的倒数和。这是降次的原因。

\medskip
\textbf{\textcolor{CExplanation}{选整体，再带范围}}
\begin{enumerate}[leftmargin=*,itemsep=0.55em]
\item 反复出现倒数和，选 $u=t+t^{-1}$，并保留 $|u|\ge2$；若 $t>0$，只保留 $u\ge2$。
\item 反复出现倒数差，选 $v=t-t^{-1}$，其范围是全体实数。
\item 只含偶次倒数和，直接选 $z=t^2+t^{-2}\ge2$，用 $t^4+t^{-4}=z^2-2$，可少一次替换。
\end{enumerate}
解方程要回代原变量；求最值要确认等号可以取得。指数式可先令 $t=a^x>0$（$a>0,a\ne1$）；两根式乘积为 $1$ 时，也可先认出倒数对再换元。
\end{explanationbox}
\clearpage
